\documentclass[11pt,a4paper]{article}
\usepackage[margin=2.4cm]{geometry}
\usepackage{booktabs,amsmath,amssymb,xcolor,hyperref,enumitem}
\usepackage[T1]{fontenc}
\hypersetup{colorlinks=true,linkcolor=blue!50!black,urlcolor=blue!50!black}
\newcommand{\have}{\textcolor{green!45!black}{\textbf{HAVE}}}
\newcommand{\todo}{\textcolor{red!70!black}{\textbf{NEEDED}}}
\newcommand{\drop}{\textcolor{gray}{\textbf{DROP}}}
\setlist{nosep,leftmargin=1.4em}
\title{\vspace{-1.5cm}\textbf{CoBit on Verifiable Tasks: a NeurIPS workshop plan}\\
\large Minimum path to submission --- deadline Sat 29 Aug (5 days)}
\author{}
\date{Mon 24 Aug 2026}
\begin{document}\maketitle\vspace{-1cm}

\section{The claim, in one paragraph}

Every published comparison on Sudoku and TinyGSM$\to$GSM8K --- ours included --- reports
\emph{one sample per problem}. Under that metric the discrete diffusion models MDLM and Duo
beat CoBit on GSM8K, but only when they use low-temperature decoding ($T{=}0.1$), which buys
pass@1 by destroying sample diversity. We argue this is the wrong measurement for
non-autoregressive reasoning models, and report the first \textbf{multi-sample (pass@$k$ /
maj@$k$)} evaluation across methods. Our thesis: \emph{CoBit trades in the opposite direction
--- it has lower pass@1 but substantially higher pass@$k$ --- and pass@$k$, not pass@1, is the
quantity that bounds what post-training can realise.}

\paragraph{Why pass@$k$ is the right proxy.} Yue et al.\ (NeurIPS 2025,
\href{https://arxiv.org/pdf/2504.13837}{arXiv:2504.13837}) show RLVR mostly \emph{reallocates}
probability mass among trajectories already in the base model's support: it raises pass@1 while
often \emph{lowering} pass@$k$ at large $k$. Base-model pass@$k$ therefore upper-bounds what RL
can extract. RL post-training for diffusion LMs exists but is masked-only
(d1 / diffu-GRPO, ICLR 2026); continuous diffusion post-training is unexplored. So
``how much headroom does a continuous bitstream dLLM have?'' is well-posed and unanswered.

\paragraph{What we will \emph{not} claim.} pass@$k$ costs $k\times$ NFE. It is a
\emph{headroom} measurement, not a matched-budget win. The paper must say this in the abstract,
or a reviewer will say it for us.

\section{Evidence already in hand (zero new compute)}

\begin{center}\small
\begin{tabular}{@{}llll@{}}
\toprule
Result & Value & Status & Use \\ \midrule
Sudoku single-sample (E/M/H) & 98.5 / 91.7 / 65.9 & \have\ published & Table 1 anchor \\
Sudoku churn sweep vs $\gamma$ & Table 16 of COLM ver. & \have & appendix \\
GSM8K single-sample, 425k & 29.19\% published & \have & Table 2 anchor \\
\quad independent reproduction & \textbf{29.42\%} (388/1319) & \have\ (this session) & validates pipeline \\
GSM8K 500k checkpoint & 29.34\% (full 1319) & \have & \emph{training is saturated} \\
Sudoku-hard maj@16 / pass@16 & $\approx$91\% / $\approx$94\% & \have\ (indicative $n$) & motivates E2 \\
GSM8K maj@32 / pass@32 & 46.5\% / 59.8\% & \have\ (256-subset) & motivates E1 \\
FKC temperature study & full negative w/ mechanism & \have & \S4 + appendix \\
\bottomrule
\end{tabular}
\end{center}

\noindent The two multi-sample rows are the reason this paper is worth writing: hard-Sudoku
maj@16 of $\approx$91\% against a published single-sample 65.9\% (Duo: 58.4\%), and GSM8K
pass@32 of 59.8\% against Duo's best-ever 36\%. Both are on reduced sample sets and must be
re-run properly --- see E1/E2 --- but the effect size is not subtle.

\section{The gap in the literature (verified, not assumed)}

I read S-FLM (\href{https://arxiv.org/abs/2605.11125}{arXiv:2605.11125}) in full. It is the
source of every baseline number in our Table 2/3.

\begin{itemize}
\item \textbf{No pass@$k$, no majority vote, no best-of-$n$, no oracle metric anywhere.}
      Sudoku: exact match, 180 steps, \emph{one sample per puzzle}. GSM8K: $T{=}1$, 1024 steps,
      \emph{one generated solution per problem}.
\item \textbf{Their checkpoints are public} (\href{https://github.com/jdeschena/s-flm}{github.com/jdeschena/s-flm},
      built on the Duo codebase). \emph{TinyGSM: all baselines released} --- AR, MDLM, Duo, FLM,
      CANDI, S-FLM. Baseline pass@$k$ therefore needs \emph{sampling only, no training}.
\item \textbf{Sudoku: no checkpoints released}, but training scripts exist and their paper
      states Sudoku trains in under 2 hours on a single L40S.
\end{itemize}

\section{Minimum experiment set (the critical path)}

Assume a 4--5 page workshop limit. Nine runs, all sampling, no training.

\begin{center}\small
\begin{tabular}{@{}llllr@{}}
\toprule
ID & Experiment & Config & Runs & GPU-h \\ \midrule
E1 & CoBit GSM8K pass@$k$ & $K{=}16$, full 1319, 1024 NFE, $\gamma{=}0.41$ & 1 & 13.5 \\
E2 & CoBit Sudoku pass@$k$ & $K{=}16$, full 2k $\times$ 3 diff., 180 NFE & 3 & $\sim$6 \\
E3 & \textbf{Duo GSM8K pass@$k$ $\times$ T} & $K{=}8$, $T\in\{1.0,\,0.1\}$, 1024 NFE & 2 & 14 \\
E4 & \textbf{MDLM GSM8K pass@$k$ $\times$ T} & $K{=}8$, $T\in\{1.0,\,0.1\}$, 1024 NFE & 2 & 14 \\
E5 & S-FLM GSM8K pass@$k$ & $K{=}8$, $T{=}1$, top-1 velocity & 1 & 7 \\
\midrule
& & & \textbf{9} & \textbf{$\sim$55} \\
\bottomrule
\end{tabular}
\end{center}

\paragraph{E3/E4 are the paper.} They test the falsifiable prediction: if Duo's pass@8
\emph{falls} as $T$ drops to the 0.1 setting that wins it pass@1, the diversity-for-headroom
trade is demonstrated rather than asserted. Everything else is context.

\paragraph{Two technical requirements.}
(i) Use the unbiased Codex estimator $\text{pass@}k = 1-\binom{n-c}{k}\big/\binom{n}{k}$, so one
$K{=}16$ run yields the \emph{whole} curve $k{=}1\ldots16$ rather than a single point; the
matched cross-method region is then $k\le 8$.
(ii) All numbers on the \emph{full} 1319 / full 2k. Our first-256 subset runs $\approx$6 points
hot (35.6\% vs 29.4\% at identical config) --- a subset would silently invalidate every
comparison.

\section{Deliberately dropped}

\begin{itemize}
\item \drop\ \textbf{Sudoku baselines.} Needs training 4 models $\times$ 3 difficulties. Cite
      published single-sample numbers instead and report CoBit pass@$k$ alone on Sudoku, stating
      plainly that no prior work reports it.
\item \drop\ \textbf{FLM / CANDI on GSM8K.} They score 0.3\% / 0.2\% single-sample; pass@$k$ of a
      model that never solves anything is uninformative.
\item \drop\ \textbf{New temperature/FKC compute.} The study is complete; it costs nothing to write up.
\item \drop\ \textbf{LM1B, OWT, CoBit-M, CFG, $k>16$.} Different paper.
\end{itemize}

\section{HPC logistics (CSD3)}

The TinyGSM checkpoints already live on CSD3 at
\texttt{/rds/project/rds-LlrDsbHU5UM/gb511/projects/BitstreamDiffusion/runs/tasks/tinygsm} ---
this is where we pulled them from, so E1 needs no data movement. Steps:

\begin{enumerate}
\item Push branch \texttt{tasks/fkc-temperature} (the GSM8K pass@$k$ eval wiring is currently
      \emph{uncommitted local work}: \texttt{gsm8k\_eval.py}, \texttt{sandbox\_gsm8k.py},
      plus the sweep scripts). \textbf{This is the one true blocker} --- nothing runs on HPC until it is pushed.
\item Clone \texttt{s-flm} on CSD3, create its env, pull the TinyGSM checkpoints from HuggingFace.
\item Add a pass@$k$ wrapper to their sampler: generate $K$ samples/problem, keep per-problem
      correctness vectors, apply the unbiased estimator. Their eval already executes programs
      in a sandbox, so this is a loop and a JSON field.
\item Submit as array jobs, one cell per (model, $T$); each cell is a few GPU-hours.
\end{enumerate}

\section{Timeline}

\begin{center}\small
\begin{tabular}{@{}lll@{}}
\toprule
Day & Work & Gate \\ \midrule
Mon 24 (today) & push branch; submit E1 + E2 on CSD3 & jobs queued \\
Tue 25 & s-flm env + checkpoints; pass@$k$ wrapper; submit E3/E4 & \textbf{go/no-go on baselines} \\
Wed 26 & E3/E4 land; E5 if clean; first figures & thesis confirmed or flipped \\
Thu 27 & write: intro, setup, results, FKC section & full draft \\
Fri 28 & polish, appendix, reproduction checks & submission-ready \\
Sat 29 & buffer + submit & --- \\
\bottomrule
\end{tabular}
\end{center}

\section{Paper skeleton}

\begin{itemize}
\item \textbf{\S1 Intro} (0.75p). Single-sample accuracy mismeasures non-AR reasoning models;
      low-$T$ buys pass@1 with diversity; pass@$k$ bounds post-training headroom.
\item \textbf{\S2 Setup} (0.5p). S-FLM parity protocol, unchanged; multi-sample extension; the
      $k\times$ NFE caveat stated up front.
\item \textbf{\S3 Results} (1.5p). \textbf{Table 1} Sudoku pass@$k$/maj@$k$ vs published single-sample.
      \textbf{Table 2} GSM8K pass@$k$ across methods $\times$ temperature.
      \textbf{Fig 1} pass@$k$ vs $k$ per method. \textbf{Fig 2} pass@8 vs $T$ for Duo/MDLM (the money plot).
\item \textbf{\S4 Why not sharpening?} (0.5p). The FKC negative, with mechanism --- free content.
\item \textbf{\S5 Discussion} (0.25p). Headroom for continuous-dLLM post-training; d1/diffu-GRPO is masked-only.
\item \textbf{Appendix.} FKC details, protocol, the 29.42\% reproduction, saturation at 500k.
\end{itemize}

\section{Risks, and what we do about each}

\begin{center}\small
\begin{tabular}{@{}p{5.3cm}p{9.3cm}@{}}
\toprule
Risk & Mitigation \\ \midrule
\textbf{Duo's pass@$k$ does \emph{not} fall at low $T$} --- thesis flips &
This is still a publishable finding (``low-$T$ is free: it costs no headroom''), but the paper
becomes framing~2, ``strongest continuous model on verifiable tasks'', with pass@$k$ as support.
Decide Wed, not Fri. \\[2pt]
s-flm env / checkpoint mismatch eats a day &
Fall back to CoBit-only pass@$k$ on both tasks + the observation that no prior work reports it.
Weaker but still a complete paper. Gate this Tue. \\[2pt]
CSD3 queue times &
Submit E1/E2 today; use short array cells rather than one long job. \\[2pt]
Reviewer: ``pass@$k$ is just more compute'' &
Address in the abstract; report matched-NFE single-sample alongside; cite Yue et al.\ for why
the quantity matters independently of budget. \\[2pt]
Config traps (\texttt{sigma\_data}, test subset) &
Both bit us this week. \texttt{sigma\_data}=0.399844765663147 must be passed explicitly (config
default 0.5 is wrong and silent); all numbers on full test sets. Re-verify every table row against
its JSON before submission. \\
\bottomrule
\end{tabular}
\end{center}

\section*{If you only do one thing}

Run \textbf{E3} (Duo pass@8 at $T{=}1.0$ vs $T{=}0.1$). It is 14 GPU-hours, needs no training,
and it alone decides whether the paper has a thesis or a table.

\end{document}
